\Askhsh{38852}{4}
\begin{ylh}{2}
    \item 3.1. Εξισώσεις 1ου Βαθμού
    \item 4.1. Ανισώσεις 1ου Βαθμού
    \item 6.1. Η Έννοια της Συνάρτησης
\end{ylh}
Ένα από τα επακόλουθα της υπερθέρμανσης του πλανήτη μας είναι το λιώσιμο των πάγων. Δώδεκα χρόνια μετά το λιώσιμο των πάγων αρχίζουν να αναπτύσσονται στους βράχους μικροσκοπικά φυτά, που ονομάζονται λειχήνες. Κάθε λειχήνα αναπτύσσεται σε σχήμα περίπου κυκλικό. Ο παρακάτω τύπος χρησιμοποιείται, για να υπολογιστεί κατά προσέγγιση η διάμετρος ($\delta$) της λειχήνας σε σχέση με την ηλικία της:

\begin{equation*}\delta(t) = 7 \cdot \sqrt{t - 12},\ \ \end{equation*}

όπου $\delta$ η διάμετρος της λειχήνας σε $mm$ και $t$ ο αριθμός των χρόνων που έχουν περάσει μετά το λιώσιμο των πάγων.

\begin{alist}
\item Να βρείτε το πεδίο ορισμού της συνάρτησης.
\monades{6}
\item Ποια είναι η διάμετρος που θα έχει μια λειχήνα 16 χρόνια μετά το λιώσιμο των πάγων;
\monades{6}
\item Η Άννα μέτρησε τη διάμετρο μιας λειχήνας που βρήκε σε κάποιο μέρος και είδε ότι ήταν $35mm$. Πόσα χρόνια έχουν περάσει από το λιώσιμο των πάγων σε αυτό το μέρος;
\monades{6}
\item Αν μια λειχήνα διπλασίασε την επιφάνειά της τα τελευταία 6 χρόνια, πριν από πόσα χρόνια δημιουργήθηκε;

\monades{7}

Δίνεται ότι η επιφάνεια κύκλου με διάμετρο δ είναι : $E = \pi \cdot \dfrac{\delta}{4}^{2}$ .

{\footnotesize Το παραπάνω θέμα αναπτύχθηκε στο πλαίσιο του έργου: «Ανάπτυξη Δοκιμασιών Αξιολόγησης Δεξιοτήτων Εγγραμματισμού στα μαθήματα της Νεοελληνικής Γλώσσας και Λογοτεχνίας, της Άλγεβρας, της Φυσικής και της Χημείας Α' Λυκείου Γενικού Λυκείου» Ανάδοχος: «Ειδικός Λογαριασμός Κονδυλίων Έρευνας (Ε.Λ.Κ.Ε) Πανεπιστημίου Ιωαννίνων» ΑΔΑΜ: 25SYMV016348911 2025-02-20.}
\end{alist}

